% !TEX root = chapter-standalone.tex

\section{Strict monoidal categories}\label{sec:stric-monoidal-cats}

\todotext{Add some introduction.
    Why do we care about this?}
    
    
    
\subsection{Formal definition}

\begin{ctdefinition}[Strict monoidal category]
    \label{def:strict-monoidal-cat}
    A \maindef{strict monoidal category} is a category $\CatC$ equipped with a \emph{strict monoidal structure}. A \emph{strict monoidal structure} is

    \constit

    \begin{enumerate}
        \item a functor $\mtimescat \colon \CatC \Ctimes \CatC \fto \CatC$, called the \emph{monoidal product}
        \item an object $\idmoncat \setin \Obof{\CatC}$, called the \emph{monoidal unit};
    \end{enumerate}

    \condit

    \begin{enumerate}
    \item the monoidal product is associative on objects: 
    \begin{equation}
     (\Obja \mtimescat \Objb) \mtimescat \Objc = \Obja \mtimescat (\Objb \mtimescat \Objc);
    \end{equation}

    \item the monoidal product is associative on morphisms,
    \begin{equation}
     (\mora \mtimescat \morb) \mtimescat \morc = \mora \mtimescat (\morb \mtimescat \morc);
    \end{equation}
    
        \item  the monoidal unit is neutral on the level of objects,
              \begin{equation}\label{eq:strict-mon-stacking-cat-cond1}
                  \Obja \mtimescatob \idmoncat = \Obja \qquad \text{and} \qquad \idmoncat \mtimescatob \Obja = \Obja;
              \end{equation}
        \item the monoidal unit is neutral on the level of morphisms,
              \begin{equation}\label{eq:strict-mon-stacking-cat-cond2}
                  \mora \mtimescatmor \catid_\idmoncat = \mora \qquad \text{and} \qquad \catid_\idmoncat \mtimescatmor  \mora = \mora.
              \end{equation}
    \end{enumerate}
\end{ctdefinition}

\todotext{Talk about the graphical notation: we can write the 1 wire as just empty spaces. (Note that so far we always had morphisms as boxes with one wire on the left and right.)}

\todotext{give example that <Set> and its cousins are strict monoidal categories}

\todotext{@JL: make/include more examples of strict monoidal categories... there are quite a few nice ones i think... scour Baez materials...}

\subsection{Non-examples}


\subsection{Basic examples}


\subsection{Applied examples}

\begin{example}
    \LTI, equipped with previously described stacking operations and an appropriate unit, is a strict monoidal stacking category.
    The unit is given by the object 0, and its identity morphism is given by the LTI system
    \begin{equation*}
        \catid_\idmoncat=\tup{\mat{0}^{0\times 1}, \mat{0}^{0\times 0}, \mat{0}^{0\times 0}, \mat{0}^{0\times 0},\mat{0}^{0\times 0}}.
    \end{equation*}
    On the side of objects, clearly $l+0=0+l=l$ for any object $l\setin \Ob_\LTI$.
    Consider $\mora\colon l\mto m$.
    On the side of morphisms we have:
    \begin{equation*}
        \begin{aligned}
            \mora \mtimescatmor \catid_\idmoncat & =
            \langle\begin{bmatrix}\prstart \\ \mat{0}^{0\times 1}\end{bmatrix},
            \begin{bmatrix}\mat{A}&\mat{0}\\ \mat{0}&\mat{0}^{0\times 0}\end{bmatrix},
            \begin{bmatrix}\mat{B}&\mat{0}\\ \mat{0}&\mat{0}^{0\times 0}\end{bmatrix},
            \begin{bmatrix}\mat{C}&\mat{0}\\ \mat{0}&\mat{0}^{0\times 0}\end{bmatrix},
            \begin{bmatrix}\mat{D}&\mat{0}\\ \mat{0}&\mat{0}^{0\times 0}\end{bmatrix}\rangle \\
                                                 & =\genericlti{}
        \end{aligned}
    \end{equation*}
    Similarly:
    \begin{equation*}
        \begin{aligned}
            \catid_\idmoncat \mtimescatmor \mora & =
            \langle\begin{bmatrix}\mat{0}^{0\times 1}\\ \prstart\end{bmatrix},
            \begin{bmatrix}\mat{0}^{0\times 0}&\mat{0}\\\mat{0}&\mat{A}\end{bmatrix},
            \begin{bmatrix}\mat{0}^{0\times 0}&\mat{0}\\\mat{0}&\mat{B}\end{bmatrix},
            \begin{bmatrix}\mat{0}^{0\times 0}&\mat{0}\\\mat{0}&\mat{C}\end{bmatrix},
            \begin{bmatrix}\mat{0}^{0\times 0}&\mat{0}\\\mat{0}&\mat{D}\end{bmatrix}\rangle \\
                                                 & =\genericlti{}
        \end{aligned}
    \end{equation*}
\end{example}



\subsection{Canonical examples}

\subsection{Constructions}


\subsection{Theory remarks}


\begin{example}
    Consider the associative stacking category from \cref{ex:assoc-stacking-semicat-integers-max}, where objects are integers and stacking of objects is taking their maximum.
    There is no possible monoidal unit here: it would have to be a neutral element for the operation ``max'', but such does not exist for $\wnumbers$.
    However, we could modify this example, and replace $\wnumbers$ with a bounded set of numbers, such as $\natnumbers$.
    Then the smallest number in $\natnumbers$, namely $0$, serves a neutral element for ``max'' and provides a monoidal unit $\idmoncat$.
\end{example}

\begin{example}
    Consider the associative stacking category from \cref{ex:assoc-stacking-semicat-lists-concat}.
    A monoidal unit would need to be a neutral element for list concatenation.
    This would be the empty list.
    In \cref{ex:assoc-stacking-semicat-lists-concat} we specified that objects are only non-empty lists, hence we do not have a strict monoidal semicategory.
    However, this example can the be easily adjusted to include also the empty list, in which case we \emph{do} obtain a strict monoidal semicategory.
\end{example}




